% !PRINTED 132
\weberpage{132}
Ist, wie vorhin, $h$ der Grad von $v$, und setzen wir, wenn $h$ durch $p$ theilbar ist, $\lambda = h$, und wenn $h$ nicht durch $p$ theilbar ist, $\lambda = hp$, so folgt aus (7) $(va)^\lambda = 1$, d. h. der Grad von $va$ ist ein Theiler von $h$ oder von $hp$, also jedenfalls ein Theiler von $p\nu$; d. h. $va$ \gesperrt{ist in} $V$ \gesperrt{enthalten}.

Ersetzen wir nun in (7) $a$ durch $a^x$ und lassen $x$ die Reihe der Zahlen $0, 1, \ldots p-1$ durchlaufen, so durchläuft $a^x$ die Gruppe $Q$ und $va^x$ das System $vQ$, und es ist für jedes $\lambda$
\weq{8}{(va^x)^\lambda = v^\lambda a^{x\lambda}.}
Wenn nun der Grad $h$ von $v$ nicht durch $p$ theilbar ist, so ist
\[
(va^x)^h = a^{xh},
\]
also nur dann $= 1$, wenn $x = 0$ ist; dagegen ist
\[
(va^x)^{ph} = 1,
\]
d. h. der Grad von $va^x$ ist, wenn $x$ nicht $= 0$ ist, ein Theiler von $ph$, aber nicht von $h$, also durch $p$ theilbar; dagegen ist er $= h$, wenn $x = 0$ ist.

Ist aber der Grad $h$ von $v$ durch $p$ theilbar, so ist er von der Form $pg$, und $g$ ist durch $p$ nicht theilbar, weil $p$ nur einfach in $n$, also auch in $h$ aufgeht. Es ist also $v^g$ ein Element aus $R$ vom Grade $p$, und muss also nach \S.\,30 in $Q$ enthalten sein. Setzen wir also hiernach
\[
v^g = a^y,
\]
so wird
\[
(va^x)^g = a^{y+gx},
\]
und dies ist dann und nur dann $= 1$, wenn $y + gx \equiv 0 \ (\mathrm{mod}\ p)$ wird, was nur für \gesperrt{einen} Werth von $x$ eintritt. Für diesen Werth von $x$ ist der Grad von $va^x$ ein Theiler von $g$; für die anderen Werthe von $x$ ist er Theiler von $pg$, aber nicht von $g$, also durch $p$ theilbar. Damit sind die beiden ersten Behauptungen des Satzes $\beta$) erwiesen.

Um auch den letzten Theil einzusehen, bemerke man, dass, wenn $v_1, v_2$ zwei Elemente aus $V$ bedeuten, die Systeme $v_1 Q, v_2 Q$ entweder ganz identisch sind, oder kein einziges gemeinschaftliches Element haben. Daraus folgt, dass man $V$ in der Weise darstellen kann:
\[
V = v_1 Q + v_2 Q + \cdots + v_{r_2} Q
\]
(obwohl $V$ im Allgemeinen keine Gruppe ist), und in jedem dieser $r_2$ Theilsysteme $vQ$ giebt es $p-1$ Elemente, deren Grad
%CONT>>
